\section{Conclusion}\label{sec:conclusion}

This project developed a modular prototype for traffic monitoring and
short-term forecasting, comprising a vehicle recognition module, a traffic
forecasting module and a web-based workbench. The system evaluates visual
recognition and sensor-based forecasting on their respective data domains while
keeping the boundary between those domains explicit.

For vehicle recognition, fine-tuning improved mAP@0.5 from 0.430 to 0.720 on the
human-labeled MIO-TCD test subset. Joint training with target-domain
pseudo-labels also improved agreement with the teacher model on the LTA
validation set relative to MIO-TCD-only training. On 63 human-labeled LTA
images, the joint model matched the pretrained detector in mAP@0.5 (0.478 versus
0.467) while reducing counting error and raising recall. These findings show the
value of traffic-specific adaptation on the benchmark, while adaptation to the
Singapore viewpoint is only partly solved; nighttime frames remain the least
reliable part of the evaluation.

Two further components extend the recognition output. A per-camera filter, calibrated by
hand, raised precision on 63 human-labeled LTA images from 0.816 to 0.833 at a
small cost in recall. A stacked ensemble of the detector count, background
differencing and zero-shot CLIP rated all 15,316 frames as free-flowing,
moderate or congested, agreeing with 74.3\% of 400 held-out human labels. The
resulting heatmaps show recurring congestion on the Causeway and Second Link
cameras at weekday peaks and on Saturdays.

For traffic forecasting, eight models were evaluated on PEMSD8 under native and
unified-input configurations. GWNET achieved the lowest MAE among the evaluated
configurations, while MultiHeadSTGCN produced competitive results without a
clearly established difference from MTGNN or AGCRN across five seeds. The staged
ablation associated direct multi-step prediction and calendar features with
substantial reductions in MAE relative to the reference STGCN configuration.
Because input features, loss and training strategy also changed, these
reductions cannot be attributed to architecture alone.

The workbench presents the monitoring and forecasting outputs in one interface,
with source labels, historical replay and a threshold-based risk indicator. The
two modules currently operate on different data domains and do not constitute an
end-to-end Singapore forecasting pipeline. Future integration requires
calibrated, time-aligned local measurements, retraining on the local series and
evaluation of the resulting workflow. Overall, the project establishes a modular
implementation and an experimental basis for that next stage.
